% !TEX root = main.tex
\section{실험 평가}
\textbf{자료 범위.} 동료의 평가 보고서를 TeX로 재구성하고 보존된 Raw와 후속 검증 결과를 반영했다.
새로운 성능 실험은 실행하지 않았다.
M1-HF, 기본 실험, 보충 실험은 서로 다른 실행이므로 수치와 통계 조건을 구분한다.
논문 초안의 두 키 구성 전체를 구현했다는 의미는 아니다.\par
\subsection{실험 환경 및 방법}
기본 환경은 표 1과 같다. 별도 설명이 없으면 Audit을 포함한 zkDPP Event를 다룬다.
\begin{center}\small\setlength{\tabcolsep}{6pt}
    \textbf{표 1. 공통 소프트웨어 설정}\par\smallskip
    \begin{tabular}{@{}ll@{}}\toprule
        \textbf{항목}        & \textbf{설정}                        \\ \midrule
        언어                 & Go 1.25.7                          \\
        Circuit            & gnark 0.15.0 / gnark-crypto 0.20.1 \\
        증명계 / 곡선           & PLONK-KZG / BLS12-381              \\
        감사 암호화             & Jubjub / Poseidon2                 \\
        병렬 설정              & GOMAXPROCS=8                       \\
        노드 / 컴파일러          & Anvil 1.7.1 / Solidity 0.8.30      \\
        EVM / Chain ID     & Prague / 31337                     \\
        Transaction gas 한도 & 30,000,000                         \\
        \bottomrule\end{tabular}\end{center}

\textbf{측정 단위와 범위}
\begin{itemize}
    \item Constraints는 컴파일한 Circuit의 제약 수다. Gas는 성공 receipt의 gasUsed다.
    \item Event E2E는 필요한 입력·path·Policy를 확보한 뒤 witness 준비부터 성공 receipt 확인까지다. 로컬 Proof Verify와 결과 대조는 타이머 밖이다.
    \item Tracing은 고유 방문 AuditRecord 수를, Freeze는 동결 대상 수를 작업량으로 사용한다. 두 작업의 시간을 더해 AuditAndFreeze E2E라고 표시하지 않는다.
\end{itemize}
\textbf{반복 조건}
\begin{itemize}
    \item 기본 Event·Audit 비교·Tracing·Freeze: warm-up 1회, 공식 5회, median.
    \item Process 보충 실험·SharedSecret 직접 전달: 조건별 공식 1회, warm-up 없음.
    \item 사전 검증 후 키 복구: warm-up 1묶음, 공식 1묶음. 묶음마다 1,000회 복구하며, 표에는 공식 묶음의 1회당 평균을 사용한다.
\end{itemize}
\textbf{환경의 적용 범위}\par
보충 Process에는 macOS 26.3, 호스트 논리 CPU 10개·메모리 32 GiB, Anvil 2 CPU·4 GiB가 기록되어 있다.
이를 기본 실험 전체의 측정 환경으로 소급하여 대입하지 않는다.
기본 실험의 하드웨어·OS는 해당 실행 근거를 별도로 명시해야 한다.
    {\small 근거: [D1]--[D6]. 시간 통계가 다른 실행을 하나의 반복 집합으로 합치지 않았다.}\par
